\section{Rational equality parameters}\label{sec:rational-parameters}

We first determine when the real incidence bound can be tight while the
limiting number of \(n\)-secants through a point of the arc is an integer.
Suppose that
\[
 \frac nq\longrightarrow\alpha,\qquad
 \frac{k}{q^2}\longrightarrow\beta,
\]
that equality holds asymptotically in \eqref{eq:classical-arc-bound}, and that
the limiting value of \(\rho_n(P)\) is an integer \(a>\lambda\).  Equality in
the incidence and coverage bounds gives
\begin{equation}\label{eq:rational-system}
 a=\frac{\lambda\alpha(1-\beta)}{(1-\alpha)\beta},
 \qquad
 \lambda(\alpha-\beta)^2=\beta(1-\alpha).
\end{equation}
If \(d=a-\lambda\), eliminating \(\beta\) gives
\begin{equation}\label{eq:rational-discriminant}
 d^2\alpha^2-d(d+1)\alpha+a=0,
 \qquad
 (d+1)^2-4a=(a-\lambda-1)^2-4\lambda.
\end{equation}

\begin{lemma}[Classification of the rational parameters]
\label{lem:rational-parameters}
\coverage{fragment}
Fix an integer \(\lambda\ge1\).  The triples
\((\alpha,\beta,a)\in\mathbb Q^2\times\mathbb Z\) satisfying
\eqref{eq:rational-system},
\[
 0<\beta<\alpha<1,\qquad a>\lambda,
\]
are in bijection with ordered factorizations \(\lambda=uv\) in positive
integers.  They are
\begin{equation}\label{eq:rational-parameter-values}
 d=u+v+1,\qquad
 \alpha=\frac{u+1}{u+v+1},\qquad
 \beta=\frac{u}{u+v+1},\qquad
 a=(u+1)(v+1).
\end{equation}
\lean{rationalEqualityParameters_satisfy_equations}
\end{lemma}

\begin{proof}
The last expression in \eqref{eq:rational-discriminant} must be an integer
square.
Factoring the difference of squares gives two positive factors of
\(4\lambda\) with the same parity.  Both factors are even; write them as
\(2u\) and \(2v\).  Then \(uv=\lambda\), and substitution gives
\eqref{eq:rational-parameter-values}.  Conversely, the displayed values
satisfy \eqref{eq:rational-system}.  The inequalities on \(\alpha\) and
\(\beta\) make the factors positive, and their order distinguishes the two
roots, so the correspondence is one-to-one.
\end{proof}

We next control \(T=\tau_n\) near one of these parameters.  The argument is
stated separately because it is also used in the modular refinements.

\begin{lemma}[Number of \(n\)-secants near equality]
\label{lem:secant-number-near-equality}
\coverage{absent}
Fix \(u,v\ge1\), put \(d=u+v+1\), and suppose
\[
 q=dm,\qquad n=(u+1)m+1,\qquad
 k=udm^2+cm+O(1),
\]
where \(c\) remains bounded.  If every external point is on at least
\(uv\) \(n\)-secants and
\begin{equation}\label{eq:F-definition}
 F(T):=\Phimin(k,nT)
 +\Phimin(q^2+q+1-k,vmT)-\binom T2
\end{equation}
is nonpositive, then
\begin{equation}\label{eq:T-localization}
 T=ud(v+1)m+O_{u,v,c}(1).
\end{equation}
\end{lemma}

\begin{proof}
Set \(b=q^2+q+1-k\).  The external incidence count gives
\begin{equation}\label{eq:T-minimum}
 T\ge T_0:=\left\lceil\frac{uvb}{vm}\right\rceil
 =ud(v+1)m+O(1).
\end{equation}
Direct substitution in \eqref{eq:phi-min} gives
\(F(T_0)=O_{u,v,c}(m)\).  Moreover,
\begin{equation}\label{eq:F-difference}
 F(T+1)-F(T)
 =\sum_{i=0}^{n-1}\left\lfloor\frac{nT+i}{k}\right\rfloor
 +\sum_{i=0}^{vm-1}\left\lfloor\frac{vmT+i}{b}\right\rfloor-T.
\end{equation}
From \(T_0\) through an interval of length \(vm+O(1)\), the first family of
floors is at least \((u+1)(v+1)-1\), and the second is at least \(uv\).
Thus \eqref{eq:F-difference} is at least
\(vm-(T-T_0)+O(1)\).  Beyond this interval, use
\(\lfloor x\rfloor\ge x-1\) and
\[
 \frac{n^2}{k}+\frac{v^2m^2}{b}-1
 =\frac1{u(v+1)}+O(m^{-1}).
\]
By the midpoint of the first interval, \(F\) has increased by a positive
quantity of order \(m^2\), whereas \(F(T_0)=O(m)\).  For larger \(T\), a
lower bound for \eqref{eq:F-difference} is
\[
 T\left(\frac{n^2}{k}+\frac{v^2m^2}{b}-1\right)-(n+vm),
\]
which is increasing because the coefficient of \(T\) is positive and is
already positive at the end of the first interval.  It follows that
\(F(T)>0\) whenever \(T-T_0\) exceeds a constant depending only on the fixed
parameters.  Therefore \(F(T)\le0\) implies
\eqref{eq:T-localization}.
\end{proof}

\begin{proof}[Proof of Theorem~\ref{thm:factor-pair}]
Lemma~\ref{lem:rational-parameters} proves the classification.  If
\((k-udm^2)/m\) exceeds a fixed constant larger than
\(c_{\mathbb Z}\), the claimed bound is automatic.
In the remaining range, Lemma~\ref{lem:secant-number-near-equality} applies
uniformly to the actual number of \(n\)-secants, because
\eqref{eq:arc-integer-condition} implies \(F(T)\le0\).  Hence
\[
 T=ud(v+1)m+h
\]
for an integer \(h\) in a fixed finite set.  Put
\(c_m=(k-udm^2)/m\).  Substitution in the external coverage inequality
\(uvb\le vmT\) gives
\begin{equation}\label{eq:coverage-linear}
 k\ge udm^2+L(h)m-O_{u,v}(1).
\end{equation}

At these parameters the balanced internal degrees are
\((u+1)(v+1)-1\) and \((u+1)(v+1)\); the balanced external degrees are
\(uv\) and \(uv+1\).  Substitution in \eqref{eq:phi-min} gives
\begin{equation}\label{eq:pair-linear}
 \frac{F(T)}m
 =\frac{(u+v)(2uv+u+v+1)}2\bigl(C(h)-c_m\bigr)+O_{u,v}(m^{-1}).
\end{equation}
The integer condition \eqref{eq:arc-integer-condition} implies \(F(T)\le0\),
so
\[
 k\ge udm^2+C(h)m-O_{u,v}(1).
\]
The affine functions \(L\) and \(C\) are decreasing and increasing,
respectively, and meet at \(h_*\).  Their maximum over the integers is
minimized at \(\lfloor h_*\rfloor\) or \(\lceil h_*\rceil\).  Since \(h\)
ranges over a fixed finite set, the bounded terms are uniform.  This proves
\eqref{eq:factor-pair-bound}.

Expanding the real bound \eqref{eq:classical-arc-bound} gives
\(c_{\mathrm{mix}}\).  Subtraction from \(c_{\mathrm{bal}}\) gives the
factorization \eqref{eq:strict-gap}.
\uses{lem:rational-parameters,lem:secant-number-near-equality,cor:n-secant-feasibility}
\end{proof}

For the first three ordered factorizations, the integer coefficients are
\begin{equation}\label{eq:first-integer-coefficients}
 (u,v)=(1,1):\ \frac{18}{5},\qquad
 (u,v)=(1,2):\ \frac92,\qquad
 (u,v)=(2,1):\ 6.
\end{equation}
The first is the ordinary-completeness family
\(q=3m\), \(n=2m+1\).  The other two have two \(n\)-secants through every
external point and limiting ratios \(n/q=1/2\) and \(3/4\).

There is also an elementary description of the parameters for which the
coverage and pair-count inequalities are simultaneously tight to first order.

\begin{proposition}[Simultaneous first-order equality]
\label{prop:simultaneous-equality}
\coverage{absent}
Put \(M_0=2uv+u+v\).  Then \(h_*\) is an integer if and only if
\begin{equation}\label{eq:simultaneous-equality-divisibility}
 M_0\mid u^2(u+1)(u^2+u+1).
\end{equation}
In particular, the case \(v=u\) has simultaneous first-order equality for
every even \(u\).
\end{proposition}

\begin{proof}
Polynomial division in \(v\), followed by multiplication by \((2u+1)^2\),
gives
\[
 (2u+1)^2u(v+1)(u^2v+u^2-v)
 =M_0Q(u,v)+u^2(u+1)(u^2+u+1)
\]
for an integer polynomial \(Q\).  Since
\(\gcd(M_0,2u+1)=1\), integrality of \(h_*\) is equivalent to
\eqref{eq:simultaneous-equality-divisibility}.  When \(v=u\),
\[
 h_*=-\frac{u(u^2+u-1)}2,
\]
which is integral for even \(u\).
\end{proof}
